\documentclass[11pt]{article}
\usepackage[a4paper,margin=25mm]{geometry}
\usepackage{amsmath,amssymb,amsthm,hyperref}
\newtheorem{theorem}{Theorem}
\newtheorem{lemma}{Lemma}
\title{Prime periods of the chocolate numbers\\MathDB 331410}
\author{Principia Math In Progress}
\date{28 September 2026}
\begin{document}
\maketitle
\begin{abstract}
For every prime $p\equiv2,3\pmod5$ with $p\ne2$, the chocolate
numbers have eventual period $p(p-1)$ modulo $p$. More precisely,
the shift $p(p-1)/2$ multiplies the sequence by the quadratic character
of $-5$, starting at index $(p+1)/2$. The proof uses integer divided
powers, so no factorial is inverted in characteristic $p$.
\end{abstract}

\section*{Task 1: the source question and its resolution}
Let $B_n$ count the ways of breaking a $2$ by $n$ chocolate bar into
unit squares, with the order of the breaks recorded. The source
\cite{chocolate} gives $B_1=1$ and, for $n\ge2$,
\begin{equation}\label{rec}
B_n=(2n-2)!+\sum_{i=1}^{n-1}\binom{2n-2}{2i-1}B_iB_{n-i}.
\end{equation}
Its Conjecture 23 (Conjecture 3 in the preprint) asks for an eventual
period divisible by $p(p-1)$ when $p\ne2,5$ and $p\not\equiv\pm1\pmod5$.
The following identifies $p(p-1)$ itself as such a period.

\begin{theorem}\label{main}
Let $p$ be an odd prime for which $5$ is a quadratic nonresidue. Put
$T=p(p-1)/2$ and $\epsilon=\left(\frac{-5}{p}\right)\in\{1,-1\}$.
For every $n\ge(p+1)/2$,
\begin{equation}\label{half}
B_{n+T}\equiv\epsilon B_n\pmod p.
\end{equation}
Consequently $B_{n+p(p-1)}\equiv B_n\pmod p$ throughout this range.
\end{theorem}
Quadratic reciprocity gives $\left(\frac5p\right)=\left(\frac p5\right)$,
so these are exactly the primes in the source question.
The theorem specifies a period, without asserting that it is the
least period. For example, the source already proves period $3$ modulo
$3$ for $n\ge2$; our formula recovers it with $T=3$.

\section*{Divided powers and the auxiliary series}
For a commutative ring $R$, let $\mathcal H(R)$ be the ring of sequences
$A=(a_0,a_1,\ldots)$, written $A=\sum_{j\ge0}a_jt^{[j]}$, with product
\[
(AB)_j=\sum_{i=0}^j\binom ji a_i b_{j-i}.
\]
Over $\mathbb Q$, $t^{[j]}$ is $t^j/j!$. The operator $D$ is the shift
$(DA)_j=a_{j+1}$. Pascal's identity proves that $D$ is a derivation.
A series with constant coefficient $1$ has a unique inverse, recursively
in the same coefficient ring.

Define integers $d_j$ by
\begin{equation}\label{drec}
d_0=1,\qquad d_1=0,\qquad
d_{j+2}=(j^2-j-1)d_j\quad(j\ge0),
\end{equation}
and put $U=\sum d_jt^{[j]}\in\mathcal H(\mathbb Z)$. All odd $d_j$ vanish.
Writing $t=t^{[1]}$, recurrence \eqref{drec} is equivalent to
\[
(1-t^2)D^2U+U=0.
\]
Indeed, the $j$th coefficient of $t^2D^2U$ is $j(j-1)d_j$.
Let $H=(1-t^2)^{-1}$; its coefficient is $j!$ for even $j$ and zero
for odd $j$. Set $Y=-(DU)U^{-1}$. Then $Y_0=0$ and
\begin{equation}\label{riccati}
DY=H+Y^2.
\end{equation}
The coefficients of \eqref{riccati} determine $Y$ uniquely from $Y_0$.
Comparing them with \eqref{rec}, including the zero even coefficients,
gives the identity
\begin{equation}\label{Y}
Y=\sum_{n\ge1}B_nt^{[2n-1]}=-(DU)U^{-1}
\quad\hbox{in }\mathcal H(\mathbb Z).
\end{equation}
This is the source's auxiliary logarithmic derivative in integral form.
All these identities may therefore be reduced coefficientwise modulo $p$.

\section*{Proof of Theorem \ref{main}}
We now work in $\mathcal H(\mathbb F_p)$. Iterating \eqref{drec} $p$ times
and using that $j,j+2,\ldots,j+2(p-1)$ run through $\mathbb F_p$ yields
\[
d_{j+2p}=d_j\prod_{r\in\mathbb F_p}(r^2-r-1).
\]
Let $\alpha,\beta\in\mathbb F_{p^2}$ be the roots of $X^2-X-1$.
The nonresidue hypothesis implies $\alpha^p=\beta$ and $\beta^p=\alpha$.
Since $\prod_{r\in\mathbb F_p}(X-r)=X^p-X$, we obtain
\[
\prod_r(r^2-r-1)
=(\alpha^p-\alpha)(\beta^p-\beta)
=-(\alpha-\beta)^2=-5.
\]
Thus, on putting $\delta=D^p$,
\begin{equation}\label{delta}
\delta^2U=-5U,\qquad
\delta^pU=(-5)^{(p-1)/2}\delta U=\epsilon\delta U.
\end{equation}
The last equality is Euler's criterion; $p\ne5$ ensures
$\epsilon^2=1$.

For every derivation in characteristic $p$, its $p$th iterate is also
a derivation: in the iterated Leibniz rule all intermediate binomial
coefficients $\binom pi$ vanish. Hence both $D^p$ and $D^{p^2}$, and
therefore
\[
L=D^{p^2}-\epsilon D^p,
\]
are derivations. Equation \eqref{delta} gives $LU=0$. As $L$ commutes
with $D$, it also gives $L(DU)=0$. Applying $L$ to $UU^{-1}=1$ gives
$L(U^{-1})=-U^{-2}LU=0$. Identity \eqref{Y} now implies $LY=0$.

Writing $y_j$ for the coefficients of $Y$, we have proved
\[
y_{r+p^2}=\epsilon y_{r+p}\qquad(r\ge0).
\]
Choose $r=2n-1-p$, which is nonnegative exactly when $n\ge(p+1)/2$.
Both subscripts are odd, and \eqref{Y} turns the last identity into
\eqref{half}. Applying it twice proves the stated period because
$\epsilon^2=1$. This proves the theorem for every admissible prime.
\qed

\section*{Proposed follow-up tasks}
\begin{enumerate}
\setcounter{enumi}{1}
\item Determine the least eventual period modulo each admissible prime,
including a criterion for when it is smaller than the period above.
\item Determine the least starting index for \eqref{half}, and characterize
the primes for which the bound $(p+1)/2$ is optimal.
\item Establish explicit eventual periods modulo $p^e$ for $e\ge2$,
including the dependence on $e$ and on the residue class of $p$.
\item For $c\in\mathbb Z$, replace \eqref{drec} by
$d_{j+2}=(j^2-j-c)d_j$ and classify the prime periods of the
coefficients of $-(DU)U^{-1}$ in terms of $1+4c$.
\end{enumerate}

\paragraph{Verification and scope.}
The proof is uniform in $p$ and $n$. The accompanying integer program
checks independent recurrence coefficients and sample congruences;
these finite checks do not replace the proof. All 16 pages of the
current preprint and all 17 pages of the final article were read.
The dated literature audit records the checked public sources and
does not claim exhaustive knowledge of all prior work. Review here
is author-agent self-review, not independent peer review.

\begin{thebibliography}{1}
\bibitem{chocolate}
C.~Ji, T.~Khovanova, R.~Park and A.~Song,
\emph{Chocolate Numbers}, Journal of Integer Sequences \textbf{19}
(2016), Article 16.1.7.
\href{https://cs.uwaterloo.ca/journals/JIS/VOL19/Khovanova/khova9.pdf}{Final article};
\href{https://arxiv.org/abs/1509.06093}{arXiv:1509.06093v1}.
\end{thebibliography}
\end{document}
